\section{Objetivo e escopo}

O HiGFlow resolve Navier--Stokes incompressível sobre malha cartesiana em
octree (HiGTree). Corpos sólidos de geometria curva não são representáveis por
essa malha sem escada de células, e é esse o problema que a fronteira imersa
resolve: o corpo é descrito por uma malha \emph{lagrangeana} de marcadores que
não coincide com a malha \emph{euleriana} do escoamento, e o acoplamento entre
as duas se faz por um campo de força.

Este relatório cobre:

\begin{itemize}
  \item a formulação implementada e as duas variantes comparadas
        (Seção~\ref{sec:formulacao});
  \item as decisões de implementação e estrutura de dados
        (Seção~\ref{sec:implementacao});
  \item a verificação dos operadores de transferência, isolada do solver
        (Seção~\ref{sec:verificacao});
  \item o experimento --- canal com obstáculo cilíndrico, \emph{benchmark}
        DFG/Schäfer--Turek 2D-1 --- e o estudo de refinamento
        (Seções~\ref{sec:experimento} e~\ref{sec:resultados}).
\end{itemize}

\paragraph{Escopo deliberadamente restrito.} As decisões abaixo foram tomadas no
início e limitam o que este trabalho afirma:

\begin{center}
\begin{tabular}{ll}
  \toprule
  corpo    & rígido e \textbf{fixo} ($U_{\mathrm{corpo}} = 0$) \\
  alcance  & apenas o solver newtoniano \\
  malha    & uniforme (sem refinamento local) \\
  \bottomrule
\end{tabular}
\end{center}

Corpo móvel, corpo deformável e interface entre fluidos \emph{não} são
extensões diretas deste trabalho; a Seção~\ref{sec:discussao} explica por quê.

\paragraph{Uma nota sobre o que este relatório registra.} Além dos resultados,
registram-se as medidas que \emph{não} discriminaram --- oráculos que passaram
com o método errado --- porque elas são a parte transferível da experiência. Um
resultado correto obtido por um teste que não podia detectar o erro é sorte, e
sorte não se reproduz.
